\section{The actual finite isotropic averaging operator}
\label{sec:finite-isotropic-average-v033}

\status{FIA1--FIA8 reconstruct the finite cancellation needed by the
original definite theta lift.  Under the five literal hypotheses in
the source's averaging theorem, they produce its vanishing sum on
every original polynomial kernel and every original adelic translate.
No additional character, finite model, or averaging hypothesis is
retained.  The hypotheses concern the supplied original real datum;
their existence and the later smooth cuspidal extension are separate
producer obligations.  These are category B proofs, with no new
published-theorem application and no new Lean replay.}

Fix the original host \(D\), with totally real field \(F\), CM field
\(E\), conjugation \(z\mapsto\bar z\), and nonzero purely imaginary
\(\xi_0=D.\mathrm{host}.\mathrm{xi0}\).  Write
\(\mathbb A_E=\mathbb A_F\otimes_F E\), with the source's
identification of its finite part with
\(E\otimes_F\mathbb A_{F,f}\).  Its integral finite ring is
\(\widehat{\mathcal O}_E\).  All inclusions, finite projections,
and residue maps below are the original
\texttt{incA}, \texttt{toEAf}, and \texttt{resS}.

Let \(S,T\) be the original Hermitian coordinate spaces, of ranks
\(n,m\), with Gram matrices \(G_S,G_T\).  The oriented space is
\[
 X=\operatorname{Mat}_{m\times n}(\mathbb A_E),\qquad
 t^\dagger=G_S^{-1}\bar t^{\,t}G_T,\qquad
 B(x,y)=\operatorname{Tr}_{\mathbb A_E/\mathbb A_F}
                  \bigl(\xi_0\operatorname{tr}(y^\dagger x)\bigr).
 \tag{FIA.1}
\]
Thus \(B=\texttt{symplO}\), whereas the source's raw Heisenberg
operator uses its half:
\[
 b=B/2,\qquad H_y f(x)=\psi(b(x,y))f(x+y).
 \tag{FIA.2}
\]
The distinction is essential: the opposite-half modulation law in
\texttt{IsPlaceModel} uses \(\psi(B(t,y))\), not
\(\psi(b(t,y))\).  We preserve both normalizations throughout.
For actual finite frames \(\beta_S,\beta_T\), put
\[
 {\cal F}(t)=\beta_S^{-1}(t^\dagger)_f\beta_T
       =\texttt{frMapAdj}(S,T,\beta_S,\beta_T,t).
 \tag{FIA.3}
\]
The subscript \(f\) includes the source's tensor-factor interchange.
Rows and columns in all matrix identities below have these orders.

% BEGIN FIA1
\paragraph{FIA1: every local character is nontrivial and has an edge.}
\label{claim:FIA1}
Suppose \(\psi\) satisfies the literal \texttt{IsCharAF}:
it is a continuous additive character of \(\mathbb A_F\), has
absolute value one, is one on the diagonal \(F\), and is not
identically one.  For every finite place \(v\), its restriction
\(\psi_v\) to the single \(F_v\)-component is nontrivial.
Moreover there is \(x_v\in F_v\) such that
\[
 \psi_v(x_v)\ne1,\qquad
 |y|_v<|x_v|_v\ \Longrightarrow\ \psi_v(y)=1.
 \tag{FIA.4}
\]
Here \(|\cdot|_v\) denotes any multiplicative normalization of
the discrete valuation; only its ordering is used.

\emph{Proof of local nontriviality.}
We give the character argument because a nontrivial character of
an arbitrary product need not be nontrivial on every factor.
On \(F_\infty\simeq\mathbb R^{[F:\mathbb Q]}\), a continuous unitary
additive character has the form
\(\exp(2\pi iL(x))\), with \(L\) real linear.
For completeness, on a real line take the continuous argument in
a small arc around \(1\).  The additive law makes its argument
additive on a smaller interval, since the integer ambiguity is
continuous and zero at the origin.  Continuity and rational
approximation make it linear there; subdivision gives the result
on the whole line.  Multiplying the coordinate restrictions proves
the assertion in finitely many dimensions.

The finite restriction of \(\psi\) kills \(N\widehat{\mathcal O}_F\)
for some positive integer \(N\).  Indeed, continuity gives a
compact open additive subgroup whose image lies in
\(\{z\in\mathbb C:|z|=1,\ \operatorname{Re}z>0\}\).
A subgroup of the circle contained in this arc is trivial:
every nonidentity point has some integral power outside the arc.
A sufficiently divisible positive integer multiple of
\(\widehat{\mathcal O}_F\) lies in this open subgroup, by the
restricted-product topology.  Rational triviality now gives
\[
            L(Nr)\in\mathbb Z\qquad(r\in\mathcal O_F).
 \tag{FIA.5}
\]
An integral basis therefore makes \(L|_F\) rational valued.
The nondegenerate rational trace pairing produces \(a\in F\)
with \(L(r)=\operatorname{Tr}_{F/\mathbb Q}(ar)\) for \(r\in F\).
One can check nondegeneracy directly: the matrix of embeddings
of a rational basis is invertible for a separable number field,
and the trace matrix is its transpose times itself.  Since an
integral basis is a real basis of \(F_\infty\), the same trace
formula holds on \(F_\infty\).

Here is a concrete standard global character.  On \(\mathbb A_\mathbb Q\)
put
\[
 \Psi_\mathbb Q((u_\infty,u_p)_p)
   =\exp\!\left(2\pi i
             \left(u_\infty-\sum_p\{u_p\}_p\right)\right),
 \qquad
 \Psi_F=\Psi_\mathbb Q\circ
                  \operatorname{Tr}_{\mathbb A_F/\mathbb A_\mathbb Q}.
 \tag{FIA.6}
\]
The \(p\)-adic fractional part is represented in
\(\mathbb Z[1/p]\cap[0,1)\); its class modulo \(\mathbb Z\) is
additive, and only finitely many terms in the sum are nonzero.
For a rational number \(q\), the difference
\(q-\sum_p\{q\}_p\) is integral.  To see this, its denominator
has no prime divisor: at every \(p\), the \(p\)-fractional part
removes the entire negative-valuation part and the other
fractional parts are \(p\)-integral.  This proves rational
triviality, continuity, and the additive law of (FIA.6).
Applying the field trace proves the corresponding assertions
for \(\Psi_F\).

The characters \(\psi\) and \(z\mapsto\Psi_F(az)\) agree at
infinity and on \(F\).  Their quotient is consequently a
continuous character of \(\mathbb A_{F,f}\) trivial on \(F\).
The diagonal \(F\) is dense in \(\mathbb A_{F,f}\).  An explicit
approximation proves this last assertion.  Given finitely many
local targets and congruence precisions, choose a positive rational
integer \(M\) that clears their denominators.  At every prime
dividing \(M\) not in the target set, impose divisibility by \(M\);
at target primes prescribe the desired residues of \(M\) times
the target, to the requested precision after division by \(M\).
The Chinese remainder theorem in \(\mathcal O_F\) gives a
numerator \(r\).  Then \(r/M\) has all required approximations
and is integral at every other place.  Thus the quotient character
is one, and
\[
                    \psi(z)=\Psi_F(az).
 \tag{FIA.7}
\]
The original nontriviality implies \(a\ne0\).
For \(v\mid p\), the local restriction is the nontrivial standard
character of \(\mathbb Q_p\) applied to
\(\operatorname{Tr}_{F_v/\mathbb Q_p}(a\,z)\).
This trace map is onto: the trace of \(1\) is the positive degree,
which is invertible in \(\mathbb Q_p\).  Multiplication by
the nonzero image of \(a\) is onto as well.  Hence \(\psi_v\)
is nontrivial.  This establishes the precise local assertion
needed here without invoking the stronger conductor-product
formula.

\emph{Proof of the edge.}
The same small-arc argument makes the kernel of \(\psi_v\)
contain \(\varpi^d\mathcal O_v\) for some integer \(d\).
The integers \(r\) for which \(\psi_v\) is one on
\(\varpi^r\mathcal O_v\) form a nonempty upper ray, with a lower
bound because \(\psi_v\) is not identically one.  Let \(d_0\)
be its least element.  Choose
\(x_v\in\varpi^{d_0-1}\mathcal O_v\) with
\(\psi_v(x_v)\ne1\).  Its valuation is exactly \(d_0-1\).
Any \(y\) with \(|y|_v<|x_v|_v\) has valuation at least \(d_0\),
so (FIA.4) follows.  In particular \(x_v\ne0\).
This proves the original \texttt{ThetaSupport.exists\_edge}.
% END FIA1

% BEGIN FIA2
\paragraph{FIA2: split charts and residue precision.}
\label{claim:FIA2}
Let \(e\) satisfy the original
\(\texttt{IsIdemAt}(e,\mathfrak P)\).  There is a finite place
\(v\) of \(F\) and a surjective ring homomorphism
\(\pi:\mathbb A_E\to F_v\) such that
\[
 \begin{gathered}
 \pi(z)=0\ \Longleftrightarrow\ ez=0,\qquad
 \pi(e)=1,\qquad e_\infty=0,\\
 \pi(\iota_F(y))=y_v,\qquad
 \operatorname{Tr}_{\mathbb A_E/\mathbb A_F}(ez)
       =\operatorname{single}_v(\pi(z)).
 \end{gathered}                                                   \tag{FIA.8}
\]
The single component on the last line has zero infinite part.
The chart factors through the original finite projection as
\(\pi=\phi\circ\texttt{toEAf}\), with
\(\phi:E\otimes_F\mathbb A_{F,f}\to F_v\).
For an integral finite element \(u\),
\[
 |\phi(u)|_v\le1,\qquad
 \operatorname{res}_{\mathfrak P}(u)=0
       \ \Longrightarrow\ |\phi(u)|_v<1.
 \tag{FIA.9}
\]

\emph{Proof.}
The literal split-idempotent conditions say that
\(e^2=e\), \(e\bar e=0\), and \(e+\bar e\) is the identity at
one finite \(F\)-place and zero elsewhere.  Multiplying the last
identity by \(e\) shows that \(e\) is zero at every other place,
including infinity.  At the indicated place the quadratic
algebra \(E\otimes_F F_v\) has the two nonzero orthogonal
idempotents \(e,\bar e\).  Its two factors have dimensions adding
to two, so each is one dimensional over \(F_v\).
Consequently this algebra is \(F_v\times F_v\), with
\(e=(1,0)\) after choosing the first factor.  Projection to this
factor is the source's chart \(\pi\); a single finite scalar
times \(e\) gives every value, proving surjectivity.
Multiplication, the kernel assertion, and the trace identity
are now coordinate identities.  In particular the source
\texttt{SplitChart.comp} is this projection: its defining
equation \(ez=\operatorname{single}_v(\pi(z))e\) determines it
uniquely.  It annihilates the infinite component and therefore
factors through the indicated finite tensor ring.

The integral ring of the split quadratic algebra is
\(\mathcal O_v\times\mathcal O_v\).  This is also the image of
the original completed integral ring: completion of the finite
\(\mathcal O_F\)-module \(\mathcal O_E\), and then projection to
the two primes over \(v\), gives precisely these two local
integer rings.  Equivalently, integral elements are characterized
by their monic equations over \(\mathcal O_v\) in each factor.
It follows that \(\phi\) sends integral elements into
\(\mathcal O_v\).

For a nonzero \(r\in\mathcal O_E\), the second clause of
\texttt{IsIdemAt} says
\[
 r\in\mathfrak P
 \ \Longleftrightarrow\
       (er^{-1})_f\notin\widehat{\mathcal O}_E
 \ \Longleftrightarrow\ |\pi(r^{-1})|_v>1
 \ \Longleftrightarrow\ |\pi(r)|_v<1.
 \tag{FIA.10}
\]
The same conclusion is true for \(r=0\).
Thus the selected local prime is exactly \(\mathfrak P\),
not its conjugate.  The kernel of the original residue map is
\(\mathfrak P\widehat{\mathcal O}_E\): the source defines it by
the quotient by this extended ideal, identifying that quotient
with \(\mathcal O_E/\mathfrak P\).
This identification follows directly by completion and the
Chinese remainder theorem.  At every other prime
\(\mathfrak P\) becomes the unit ideal, while at \(\mathfrak P\)
completion preserves the residue field.
Every element of the extended ideal is a finite sum of
integral multiples of elements of \(\mathfrak P\).
Equation (FIA.10), integrality of \(\phi\), and the
nonarchimedean triangle inequality prove (FIA.9).
In particular
\[
       \psi\!\left(\operatorname{Tr}(ez)\right)=\psi_v(\pi(z)).
 \tag{FIA.11}
\]
No rescaling of the character, trace, or residue homomorphism
has occurred.
% END FIA2

% BEGIN FIA3
\paragraph{FIA3: the constant phase on the original top block.}
\label{claim:FIA3}
Assume \(m\le n\), \(\texttt{IsCharAF}(\psi)\), and
\(\texttt{IsIdemAt}(e,\mathfrak P)\).
Let \(c\in\mathbb A_E^m\) with \(ec_j\ne0\) for some \(j\).
The frames \(\beta_S,\beta_T\) are the actual invertible finite
matrices; no change of them is made.  Put
\[
                         q_{\mathfrak P}
                            =\#(\mathcal O_E/\mathfrak P)>0.
 \tag{FIA.12}
\]
There exist \(l\in\mathbb A_E^n\) and \(\zeta\in\mathbb C\)
such that
\[
 el=l,\qquad \zeta\ne1,\qquad \zeta^{q_{\mathfrak P}}=1,
 \tag{FIA.13}
\]
and, for every \(t\in X\),
\[
 \begin{split}
 {\cal F}(t)\in\operatorname{Mat}_{n\times m}
                      (\widehat{\mathcal O}_E),\qquad
 \operatorname{top}_m(\operatorname{res}_{\mathfrak P}{\cal F}(t))=I_m\\
 \Longrightarrow\quad
            \psi\bigl(B(t,c\,l^t)\bigr)=\zeta .
 \end{split}                                                       \tag{FIA.14}
\]
This proves \texttt{ThetaSupport.exists\_const}, with its positive
period specified explicitly.

\emph{Proof.}
Take the chart \(\pi,\phi\) of FIA2 and the edge \(x=x_v\)
of FIA1.  Define
\[
 d=\phi(\beta_T^{-1})\,\pi(c),\qquad
 u=\pi(\xi_0).
 \tag{FIA.15}
\]
The vector \(\pi(c)\) is nonzero, and the finite matrix
\(\phi(\beta_T)\) is invertible, so \(d\ne0\).
Choose \(j<m\) with \(|d_j|_v\) maximal.  It is nonzero.
The composite \(E\to\mathbb A_E\to F_v\) is a unital
field homomorphism and hence injective, so \(u\ne0\).
Set
\[
 s=-\frac{x}{u\,d_j},\qquad
 \pi(l_{0,a})=s\,\phi((\beta_S^{-1})_{ja}),\qquad
 l_a=e\,l_{0,a},\qquad \zeta=\psi_v(x).
 \tag{FIA.16}
\]
Surjectivity of \(\pi\) supplies every \(l_{0,a}\);
the row \(j\) is valid in the \(n\)-row frame because \(m\le n\).
These choices depend on \(c\), the selected place, character,
and frames, and not on \(t\).

Let \(U={\cal F}(t)\) satisfy the two conditions in (FIA.14).
Its \(j\)-th row, after applying \(\phi\), satisfies
\[
          |\phi(U_{ja})-\delta_{ja}|_v<1
                             \qquad(0\le a<m),                  \tag{FIA.17}
\]
by (FIA.9).  Put \(z=\sum_{a<m}\phi(U_{ja})d_a\).
Every term of
\[
                  z-d_j=\sum_{a<m}
                         (\phi(U_{ja})-\delta_{ja})d_a
 \tag{FIA.18}
\]
has absolute value strictly smaller than \(|d_j|_v\);
the finite ultrametric triangle inequality therefore gives
\(|z-d_j|_v<|d_j|_v\).
The actual matrix identity
\[
       \phi(U)=\phi(\beta_S^{-1})\,\pi(t^\dagger)\,
                    \phi(\beta_T)
 \tag{FIA.19}
\]
and (FIA.15)--(FIA.16) yield
\(\pi(l_0^t t^\dagger c)=s z\).

We check the orientation and sign directly.
From (FIA.1),
\(\operatorname{tr}(t^\dagger c\,l_0^t)
=l_0^t t^\dagger c\).
Hermitian symmetry, \(\bar\xi_0=-\xi_0\), and invariance
of the quadratic trace under conjugation give
\[
 B(t,c(e l_0)^t)
   =\operatorname{Tr}
          \bigl(e(-\xi_0\,l_0^t t^\dagger c)\bigr).
 \tag{FIA.20}
\]
Applying (FIA.11), its character value is
\(\psi_v(-u\,s z)\).  The chosen negative sign in \(s\)
is exactly what makes
\[
             -u\,s z=x+x\frac{z-d_j}{d_j}.
 \tag{FIA.21}
\]
The error term has absolute value \(<|x|_v\), so FIA1
annihilates it.  Additivity proves (FIA.14).

Finally \(q_{\mathfrak P}\in\mathfrak P\): the additive group
of the finite residue field is killed by its cardinality.
Consequently \(|q_{\mathfrak P}|_v<1\) by (FIA.10), and
\[
 \zeta^{q_{\mathfrak P}}
   =\psi_v(q_{\mathfrak P}x)=1
 \tag{FIA.22}
\]
by the same edge inequality.  Nontriviality of \(\zeta\)
was built into the choice of \(x\).
Its order need only divide \(q_{\mathfrak P}\); no primitive-root
claim or new choice of residue character is required.
% END FIA3

% BEGIN FIA4
\paragraph{FIA4: the actual finite half-datum meets the model's guards.}
\label{claim:FIA4}
Let \(\Lambda_S,\Lambda_T\) be full
\(\mathcal O_E\)-lattices.  In the application they are the
original self-dual lattices.  Write
\[
 \Lambda=\texttt{latO}(S,T,\Lambda_S,\Lambda_T,\mathrm{True}).
 \tag{FIA.23}
\]
It contains the actual fractional-ideal factor
\(\mathfrak A=D.\mathrm{host}.\mathrm{AA}\).  In the self-dual
case it is
\[
 \widehat{\mathfrak A}\,
 \operatorname{Hom}_{\widehat{\mathcal O}_E}
  (\widehat\Lambda_S,\widehat\Lambda_T).
\]
Replacing it by the unweighted Hom lattice would be incorrect.

Let \(P_g,Q\) be the finite sets in
\texttt{IsThPVecOf}, and let \(e_\Sigma\) be its supplied
sum of split idempotents over \(P_g\cup Q\).
Set \(a=\bar e_\Sigma\).  The actual finite function is
\[
 \Phi(t)=
 \begin{cases}
 \Omega_K(P_g,\operatorname{res}{\cal F}(t))\,
 \Omega_Q(Q,m,\chi,-1,\operatorname{res}{\cal F}(t)),
       &at=t,\ {\cal F}(t)\text{ integral},\\
 0,&\text{otherwise}.
 \end{cases}                                                     \tag{FIA.24}
\]
Here \(\chi=\texttt{xiR}\) is the stored family of residue
characters.  The first factor is one exactly when, at each
\(\mathfrak P\in P_g\), the upper \(m\times m\) block is the
identity and all rows below it are zero; otherwise it is zero.
The second factor is zero unless its indicated lower rows
vanish and top determinants at \(Q\) are units, and otherwise
is the product of the stored inverse determinant characters.
Only its dependence on the residues at \(Q\) is used.

Then \(\Phi\) satisfies the exact \texttt{IsSBHalf} guard:
its support lies in \(aX\), and there is an integer \(K>0\)
such that
\[
 \begin{gathered}
 au=u,\ u\in\Lambda
       \quad\Longrightarrow\quad
       \Phi(t+Ku)=\Phi(t)\quad(t\in X),\\
 \Phi(t)\ne0\quad\Longrightarrow\quad Kt\in\Lambda .
 \end{gathered}                                                     \tag{FIA.25}
\]
\emph{Proof.}
We first justify uniform lattice denominators.
The source's lattice \(\Lambda\) is generated over the finite
integral \(F\)-adeles by the matrices
\[
           A\,v\,\bar w^{\,t}G_S,\qquad
 A\in\mathfrak A,\quad v\in\Lambda_T,\quad w\in\Lambda_S.
 \tag{FIA.26}
\]
The fractional ideal and both lattices are finitely generated
over their integer rings.  Their rational spans are \(E\)
and the two full coordinate spaces.  Since \(G_S\) is invertible,
these matrices span \(\operatorname{Mat}_{m\times n}(E)\).
Take finite \(\mathcal O_E\)-generating families in (FIA.26)
and multiply the resulting matrix generators by an integral
\(\mathbb Z\)-basis of \(\mathcal O_E\).
Their \(\mathbb Z\)-span is a full finitely generated
lattice.  Clearing the coefficients of its generators and of
expressions for the rational matrix units gives positive integer
inclusions in both directions between \(\Lambda\) and the
standard integral finite matrix lattice.
This also proves that every element of \(\Lambda\) has zero
infinite part.

On the finite matrix space, \({\cal F}\) is the composition of
conjugate transpose with left and right multiplication by
invertible finite-adelic matrices.  Its inverse is of the same
form.  Entries of those matrices and their inverses are integral
away from finitely many places; one positive rational integer
clears their finitely many remaining denominators.
Combined with the preceding lattice inclusions, this gives
positive integers \(M,J\) such that
\[
 M{\cal F}(\Lambda)\subset
                   \operatorname{Mat}(\widehat{\mathcal O}_E),
 \qquad
 t_\infty=0,\ {\cal F}(t)\text{ integral}
                       \ \Longrightarrow\ Jt\in\Lambda .
 \tag{FIA.27}
\]
This is the content of the source's
\texttt{frMapAdj\_up/down} and lattice-denominator lemmas,
proved here for the actual fractional lattice.

Put \(q=\prod_{\mathfrak P\in P_g\cup Q}q_{\mathfrak P}\)
and \(K=JMq>0\); the empty product is \(1\).
For \(u\in\Lambda\), the increment
\({\cal F}(Ku)=Jq\,M{\cal F}(u)\) is integral and is zero
modulo every \(\mathfrak P\in P_g\cup Q\).
Thus integrality of \({\cal F}(t)\) is unchanged by this
increment, all relevant residue matrices are unchanged,
and the half condition is unchanged when \(au=u\).
Every branch of (FIA.24), including its zero branch,
is therefore invariant.  If \(\Phi(t)\ne0\), then \(at=t\).
Since \(a_\infty=0\), we have \(t_\infty=0\); the second
inclusion in (FIA.27) gives \(Jt\in\Lambda\), and hence
\(Kt\in\Lambda\).  This proves all guards in (FIA.25).
% END FIA4

% BEGIN FIA5
\paragraph{FIA5: a genuine eigencharacter of the stored finite vector.}
\label{claim:FIA5}
Suppose \(m\le n\), the original lattices are self-dual,
\(\mathfrak P\in P_g\), and
\[
 \texttt{IsThPVecOf}
   (\psi,\theta,S,T,s_S,s_T,\Lambda_S,\Lambda_T,
         \beta_S,\beta_T,P_g,Q,\chi,f_0,f)
 \tag{FIA.28}
\]
holds for the actual vector \(f\).
Assume \(c\in\mathbb A_E^m\) has nonzero projection at every
split idempotent, exactly as in \texttt{exists\_eigen}.
There exist \(l,\zeta\) with \(\zeta\ne1\),
\(\zeta^{q_{\mathfrak P}}=1\), such that, for all integers
\(k\ge0\),
\[
        (c(kl)^t)_\infty=0,\qquad
        H_{c(kl)^t}f=\zeta^k f .
 \tag{FIA.29}
\]

\emph{Proof.}
Unpack (FIA.28).  It contains a standard vacuum \(f_0\),
a supplied sum \(e_\Sigma=\sum_{\mathfrak R\in P_g\cup Q}
e_{\mathfrak R}\), and the exact
\(\texttt{IsThetaAt}\) assertion for the half \(a=\bar e_\Sigma\)
and finite function (FIA.24).
The single-prime idempotents are pairwise orthogonal.
Indeed FIA2 identifies each with the characteristic idempotent
of its \(\mathfrak R\)-component; distinct primes give disjoint
components.  Thus
\[
            e_{\mathfrak P}e_\Sigma=e_{\mathfrak P}.
 \tag{FIA.30}
\]
The vacuum's \texttt{IsStdChar} field implies the precise
\texttt{IsCharAF} hypothesis used in FIA1 and FIA3.
Apply FIA3 at \(e_{\mathfrak P}\) to obtain \(l,\zeta\).
It has \(e_{\mathfrak P}l=l\), so (FIA.30) gives
\(e_\Sigma c(kl)^t=c(kl)^t\).  This proves its zero
infinite component and makes it an element of the half
opposite to \(a=\bar e_\Sigma\).

The literal \texttt{IsThetaAt} is an \texttt{AllSome}:
at least one linear model \(T_a\) satisfies
\texttt{IsPlaceModel}, and \emph{every} such model sends
the finite function \(\Phi\) to the same supplied vector \(f\).
Choose one of those models.  FIA4 has verified the guard
needed for its modulation axiom, which gives
\[
 H_{c(kl)^t}f
   =T_a\!\left(t\longmapsto
           \psi(B(t,c(kl)^t))\,\Phi(t)\right).
 \tag{FIA.31}
\]
For \(\Phi(t)\ne0\), its \(\Omega_K\) factor forces the
top identity block at \(\mathfrak P\); its integrality guard
also holds.  FIA3 and additivity of \(B\) in its second
argument show that the character in (FIA.31) is \(\zeta^k\).
For \(\Phi(t)=0\), the equality of the two multiplied
functions holds automatically.  Linearity of the
\emph{chosen actual} model therefore gives
\(H_{c(kl)^t}f=T_a(\zeta^k\Phi)=\zeta^k f\).
This is an equality of the stored functions, not a replacement
of \(f\) by a newly normalized vector.  Neither a scalar
uniqueness theorem for Weil representations nor the existence
of a new model is needed.
% END FIA5

% BEGIN FIA6
\paragraph{FIA6: the original good-prime set is nonempty.}
\label{claim:FIA6}
If \(L:\texttt{LevCh}(D,T_0)\) and \(L.\texttt{Is}(\mu_0)\),
then \(L.P_g\ne\varnothing\).  Moreover the supplied
\(L.\beta_V(i)\), \(L.\beta_W(i)\) have the full self-dual
lattices required by FIA4--FIA5.

\emph{Proof.}
The original character frame contains
\(\texttt{PiBar.pb}(0)=\mathfrak p_0\), a nonzero prime of
\(\mathcal O_{K_0}\), where
\(K_0=D.\mathrm{host}.K\subset E\).
It also records that the positive rational prime
\(D.\mathrm{host}.g(0)\) belongs to \(\mathfrak p_0\).
There is a prime \(\mathfrak P\) of \(\mathcal O_E\) over
\(\mathfrak p_0\).  Here is a short proof of this particular
lying-over assertion.  Localize the finite torsion-free
\(\mathcal O_{K_0}\)-module \(\mathcal O_E\) at
\(\mathfrak p_0\).  Over the resulting discrete valuation
ring it is free of positive rank.  Its reduction modulo
the maximal ideal is consequently a nonzero finite-dimensional
algebra over \(k(\mathfrak p_0)\).  A maximal ideal of that
algebra pulls back to a prime of the localization and then
of \(\mathcal O_E\); its contraction is \(\mathfrak p_0\),
because the residue-field scalar map is injective.
The pulled-back prime is nonzero, as it contains
\(D.\mathrm{host}.g(0)\ne0\); in the Dedekind ring
\(\mathcal O_E\) it is a height-one prime.

By definition this says \(L.C.P.\texttt{Div}(0,\mathfrak P)\).
The literal field \(L.\texttt{Is}.P_g\) is the equivalence
\[
       \mathfrak P\in L.P_g
          \quad\Longleftrightarrow\quad
       \exists k\in\{0,1\},\ L.C.P.\texttt{Div}(k,\mathfrak P).
 \tag{FIA.32}
\]
It proves \(\mathfrak P\in L.P_g\).
This independent argument reaches the exact
\texttt{exists\_mem\_Pg} conclusion and avoids the source's
longer Galois-conjugacy detour.
Splitting at the selected \(\mathfrak P\) is supplied when
the actual \texttt{IsThPVecOf} is unpacked; it has not
been inferred from lying over alone.

Finally, the two level fields give
\(\texttt{IsSDFrame}(V,\beta,\Lambda,1)\).
This means \(\texttt{IsBaseFrame}(V,\beta,\Lambda,1^c)\)
for some \(c\).  Its two lattice inclusions identify
\(\Lambda\) with the coordinates of the original full
standard lattice \(V.L0\), and its duality field is
\(\Lambda^\vee=\Lambda\).  Hence \(\Lambda\) is finitely
generated, spans \(V\), and has precisely the
\texttt{IsSelfDualC} property used in FIA5.
All these facts concern the stored lattice and frame.
% END FIA6

% BEGIN FIA7
\paragraph{FIA7: finite Heisenberg shifts commute with the original creation words.}
\label{claim:FIA7}
Suppose \(\theta\) is the supplied \texttt{IsInfFam}, and
\(u\in X\) has zero infinite part.  For the original
\texttt{pureTensorO} and any polynomial \(P\),
\[
 \operatorname{pureTensorO}(H_u f)(P)
       =H_u\bigl(\operatorname{pureTensorO}(f)(P)\bigr).
 \tag{FIA.33}
\]
All Cartan involutions, bases, and scaling factors in this
formula are the same on the two sides.  The operation is also
homogeneous in its seed:
\[
       \operatorname{pureTensorO}(z f)(P)
          =z\,\operatorname{pureTensorO}(f)(P)
                         \qquad(z\in\mathbb C).
 \tag{FIA.34}
\]
These statements apply to every actual polynomial kernel.

\emph{Proof.}
Let \(v(r)=\texttt{lineAt}(\theta,\tau,Z)(r)\) be any
one of the real lines used to define \(\texttt{dRho}\).
The defining property of \(\texttt{IsInfFam}\) identifies
the infinite tensor algebra with its archimedean components;
the line has prescribed component \(rZ\) at \(\tau\), zero
at the other infinite places, and zero finite part.
Its existence and uniqueness follow by taking these
components and the zero finite component.  Thus it is the
source's chosen \texttt{lineAt}.
Since finite and infinite adele components multiply to zero,
\[
                B(u,v(r))=B(v(r),u)=0.
 \tag{FIA.35}
\]
The explicit formula (FIA.2) gives the Heisenberg composition
law
\[
           H_xH_y=\psi(b(x,y))H_{x+y}.
 \tag{FIA.36}
\]
Equation (FIA.35) therefore gives
\(H_{v(r)}H_u=H_uH_{v(r)}\).
For a fixed evaluation point \(w\), the right side is
\[
     \psi(b(w,u))\,(H_{v(r)}f)(w+u),
 \tag{FIA.37}
\]
whose prefactor is constant in the real variable \(r\).
Differentiating at \(r=0\) proves
\[
       \mathrm d\rho(Z)\,H_u f
           =H_u\,\mathrm d\rho(Z)f.
 \tag{FIA.38}
\]
This is pointwise differentiation of an explicit formula.
It does not pass a derivative through an arbitrary endomorphism
of the space of all functions.
It also agrees with the source's total \texttt{deriv} convention:
the phase in (FIA.37) is nonzero, so differentiability of one
side is equivalent to that of the other, and the zero value
assigned at a nondifferentiable point respects the equality.

Each creation operator is the prescribed scalar multiple of
\(\mathrm d\rho(Z)-i\,\mathrm d\rho(JZ)\).
Equation (FIA.38) and linearity of \(H_u\) give commutation
with each such operator.  Induction through the exact ordered
list defining \texttt{opPow} gives commutation for every
monomial word.  Taking the finite linear combination of
monomial words defining the polynomial proves (FIA.33);
no reordering of creation operators is necessary.
Constant multiplication in the same one-variable derivatives
proves (FIA.34), first for each word and then for its finite
polynomial sum.
For the actual seed, the given
\texttt{thP\_pt\_theta} at \(P=1\) also identifies it as
the original theta function, since
\(\operatorname{pureTensorO}(f)(1)=f\).
% END FIA7

% BEGIN FIA8
\paragraph{FIA8: the uniform original averaging operator.}
\label{claim:FIA8}
Let \(R:\texttt{RealDat}(D,A,K_s)\) be the actual supplied
datum.  Assume exactly
\[
 \begin{gathered}
 R.\texttt{toLevCh.Is}(A.\mu_0),\qquad R.\texttt{models\_inf},
 \qquad R.\texttt{thP\_}\omega_P,\\
 R.\texttt{thP\_vP},\qquad R.\texttt{thP\_pt\_theta}.
 \end{gathered}                                                    \tag{FIA.39}
\]
There is a single integer \(N_R>0\) such that for every
original index \(i\), every tuple \(t\), every
\(h\in G(V_c(i))(\mathbb A_F)\),
\(g\in G(W_c(i))(\mathbb A_F)\), and every nonzero
\(\xi\in E^{\operatorname{rank}W_c(i)}\), there is a list
\[
       a_0,\ldots,a_{N_R-1}\in
                   \mathbb A_E^{\operatorname{rank}V_c(i)}
 \tag{FIA.40}
\]
such that, simultaneously for every polynomial \(P\),
\[
 \sum_{k=0}^{N_R-1}
 H_{\operatorname{along}(\xi,a_k)}
       \left(R.\omega_P(i,t)(h,g)\,
                        R.\mathrm{ptP}(i,t)(P)\right)=0 .
 \tag{FIA.41}
\]
The equality is an equality of functions on the actual
oriented matrix space.  Here
\(\operatorname{along}(\xi,a)=\iota_E(\xi)a^t\)
is literally the source's \texttt{along}.

\emph{Proof.}
FIA6 chooses one \(\mathfrak P\in R.P_g\), once for the
whole datum; set \(N_R=q_{\mathfrak P}\).
Fix \(i,t,h,g,\xi\ne0\), and abbreviate
\(S=V_c(i)\), \(T=W_c(i)\).
The original rank formulas give
\(\operatorname{rank}T=r\le\operatorname{rank}S\);
the required self-dual lattice statements were checked in FIA6.
Let \(M_S(h),M_T(g)\) denote the actual coordinate matrices
\texttt{mt}.  Put
\[
                   c=M_T(g^{-1})\,\iota_E(\xi).
 \tag{FIA.42}
\]
For every split chart \(\pi\), the matrix \(\pi(M_T(g))\)
is invertible, while \(\pi(\iota_E(\xi))\ne0\):
the composite \(E\to F_v\) is injective.
Therefore \(\pi(c)\ne0\), exactly the nonvanishing condition
in FIA5.
Apply that claim to the \emph{stored}
\(f=R.vP(i,t)\) and its supplied \texttt{IsThPVecOf}
to obtain \(l,\zeta\) with period \(N_R\).

Define the averaging row vector by the literal formula
\[
                         a_k^t=(kl)^t M_S(h^{-1})
                          \qquad(0\le k<N_R).
 \tag{FIA.43}
\]
Its construction precedes the choice of a polynomial.
Matrix multiplication gives
\[
 \begin{aligned}
 \operatorname{along}(\xi,a_k)
   &=M_T(g)\,c(kl)^tM_S(h^{-1})\\
   &=\operatorname{actO}(S,T,h,g)\bigl(c(kl)^t\bigr).
 \end{aligned}                                                     \tag{FIA.44}
\]
By FIA7 and (FIA.29),
\[
     H_{c(kl)^t}\,R.\mathrm{ptP}(i,t)(P)
                =\zeta^k R.\mathrm{ptP}(i,t)(P).
 \tag{FIA.45}
\]
The supplied \texttt{thP\_pt\_theta} makes each polynomial
vector a theta function, so the \emph{original action's}
\texttt{IsWeilDatumO.om\_heis} applies to precisely that
vector.  Its covariance and linearity, with (FIA.44), yield
\[
 \begin{split}
 H_{\operatorname{along}(\xi,a_k)}
   \bigl(R.\omega_P(i,t)(h,g)\,R.\mathrm{ptP}(i,t)(P)\bigr)
   =\zeta^k
       R.\omega_P(i,t)(h,g)\,R.\mathrm{ptP}(i,t)(P).
 \end{split}                                                       \tag{FIA.46}
\]
Summing this finite identity proves (FIA.41), because
\[
       \sum_{k=0}^{N_R-1}\zeta^k
                   =\frac{\zeta^{N_R}-1}{\zeta-1}=0 .
 \tag{FIA.47}
\]
There is no convergence or exchange-of-integrals issue in this
step.

The source's \(\texttt{IsoVec}(T,\xi)\) is the conjunction
of \(\xi\ne0\) with its Hermitian norm being zero.
Only its first clause has been used.  Restricting (FIA.41)
to such vectors proves the exact
\texttt{ThetaSupport.\allowbreak sum\_heisO\_\allowbreak along\_eq\_zero} and its
wrapper
\texttt{fock\_vectors\_\allowbreak killed\_\allowbreak
heisenberg\_\allowbreak average\_\allowbreak isotropic}.
Our quantifiers are slightly stronger: \(N_R\) is independent
of \(i,t,h,g,\xi,P\), and the list is independent of \(P\).
No independence of the list from \(i,t,h,g,\xi\) is asserted.
The case of a zero-dimensional target has no nonzero \(\xi\),
so it causes no exception.
% END FIA8

\paragraph{What has been supplied to the cuspidal-extension branch.}
FIA8 closes the finite \texttt{hsupp} premise on every original
tuple, including the base tuple.  It uses the original level
data, actual half-datum, original creation operators and original
Weil covariance, all coupled as in (FIA.39).
The same nonzero field embedding and the same residue kernel
are used in the frame computation and the transported-vector
computation.  The construction contains no unconstrained
auxiliary root of unity, replacement of a stored vector, or
new theta action.  BTH13--BTH15's scalar-normalization argument
is compatible with this calculation but is not a dependency:
only the already supplied covariance field is used.

The constant-term invariance under \(\operatorname{along}\),
its measure normalization, and the implication from those
constant terms to cuspidality are the next separate arguments.
The original smooth multiplicity-space extension still also
requires its analytic and representation-theoretic producers.
FIA8 makes no assertion of moderate growth, lowering
annihilation, bounded holomorphic realization, simultaneous
\texttt{ioP} existence, or original geometric/rational reading.
In particular a published rapid-decay theorem is neither used
nor certified by this component.
